\subsection{F9: what the entire \'etale topos reconstructs}

In this subsection a morphism \emph{over} $S_{\mathrm{et}}$ includes a
specified invertible $2$-cell to the structure morphism. Isomorphisms
between such morphisms must respect that $2$-cell. This is essential,
including over a finite field. A geometric morphism is called
\emph{pinned} when its map on underlying spaces sends closed points to
closed points, and \emph{chivalrous} when it sends constructible subsets
to constructible subsets. These are conditions on a geometric morphism,
not additional structure attached to a scheme.

\begin{result}{F9.1}{Proved}{The 1-faithfulness assertion over the stated base}
\textup{[F9; \S2, manuscript p.\ 25; scan 52.]}
For a locally noetherian scheme $S$, a reduced $S$-scheme $X$ locally
of finite presentation, and an $S$-scheme $Y$ locally of finite
presentation, the map
\[
 \operatorname{Hom}_S(X,Y)\longrightarrow
 \pi_0\operatorname{Hom}_{S_{\mathrm{et}}}(X_{\mathrm{et}},Y_{\mathrm{et}})
 \tag{F.1}
\]
is injective. Thus the manuscript's 1-faithfulness statement does not
need its tentative characteristic-zero restriction.
\end{result}

\begin{proof}
The finite-presentation, quasi-compact quasi-separated-base assertion is
Carlson--Stix, Lemma 5.4, in version 2 (20 September 2026)
\cite{CF_CarlsonStix2026}. Its source hypothesis is that $X$ is reduced;
no assumption that its fibres are reduced is present in that lemma.
Scheme morphisms of finite presentation induce chivalrous geometric
morphisms, and the lemma gives their injectivity into the indicated
groupoid of such morphisms. Inclusion of this full subgroupoid in all
geometric morphisms does not change whether its two objects are isomorphic.

Here is the localization giving precisely the manuscript's scope. If
$f_{\mathrm{et}}\simeq g_{\mathrm{et}}$, their maps on underlying
spaces coincide, because open subsets are the subterminal objects of
the \'etale topos. Choose an affine open in $S$, an affine open in $Y$
over it, and affine opens in its common inverse image under $f,g$.
On each such source open, both restricted maps have affine target
and are between finite-presentation schemes over the chosen affine
base. Apply the cited lemma there. These source opens cover $X$, so
$f=g$. This proves the reduction even when $S,X,Y$ are not
quasi-compact.

Absolute Frobenius is not a counterexample to (F.1). Over
$\mathbb F_p$, the natural isomorphism on \'etale covers induced by
relative Frobenius restricts on the base to the arithmetic Frobenius
automorphism of the identity functor of $B G_{\mathbb F_p}$. That
automorphism is nonidentity. Thus it does not give the required
isomorphism over $S_{\mathrm{et}}$ with its fixed $2$-cell.
\end{proof}

\begin{result}{F9.2}{Known}{The precise full-faithfulness theorems}
\textup{[F9; \S2, manuscript pp.\ 25--26; scans 51--52.]}
Let $\mathrm{UH}$ denote the universal homeomorphisms. The functor
\[
 (\operatorname{Sch}^{\mathrm{ft}}_k)[\mathrm{UH}^{-1}]
 \longrightarrow \operatorname{RTop}^{\mathrm{pin}}_{k_{\mathrm{et}}},
 \qquad X\longmapsto X_{\mathrm{et}},                         \tag{F.2}
\]
is fully faithful for every infinite field $k$ finitely generated over
its prime field, and for every sub-$p$-adic field $k$ (a subfield of a
finitely generated extension of some $\mathbb Q_p$). The first assertion
is Carlson--Haine--Wolf, Theorem 0.1; the second is Carlson--Stix,
Theorem A (=Theorem 4.1). In characteristic zero, the localization is
represented by seminormalization; in positive characteristic it is
represented by perfection.

There is also the following relative assertion. Let $S$ be noetherian,
Nagata, geometrically unibranch, and suppose each generic residue field
of $S$ is sub-$p$-adic. Restrict to finite-type $S$-schemes each of whose
irreducible components is universally open over $S$. After inverting
$\mathrm{UH}$, the topos functor into chivalrous geometric morphisms
over $S_{\mathrm{et}}$ is fully faithful. This is Carlson--Stix,
Theorem D (=Theorem 6.12 and Proposition 6.11).
\end{result}

\begin{proof}
These are the exact statements of the opened primary sources
\cite{CF_CHW2024,CF_CarlsonStix2026}. To derive the assertions about
the manuscript's original objects, observe first that a universal
homeomorphism induces an equivalence of \'etale topoi, by topological
invariance of the \'etale site. Thus the functor factors through the
displayed localization. In characteristic zero the universal
homeomorphism $X^{\mathrm{sn}}\to X$ represents the localization:
maps in it can be calculated from the seminormalization of the source.
Accordingly, for a seminormal finite-type source over the fields in
(F.2), the comparison on actual scheme Hom-sets is already a bijection;
for two seminormal schemes, topos equivalences are exactly scheme
isomorphisms. Over the relative bases just stated, normal sources are
absolutely weakly normal and the analogous actual-Hom comparison
holds; this is Corollary E of Carlson--Stix. These conclusions follow
by applying the displayed fully faithful functors to their localized
objects, not by identifying the original schemes without their
normality hypothesis.

For a general base, chivalrous cannot be replaced by pinned in the
relative statement. Take a DVR $R$, an element $a$ of positive
valuation, and the closed point $(aT-1)$ of $\mathbb A^1_R$.
The projection $\mathbb A^1_R\to\operatorname{Spec}R$ sends it to the
generic point. Its topos morphism is therefore not pinned, although
it is induced by a scheme morphism in the stated category. Its image
of any constructible set is constructible by Chevalley's theorem.
This also gives the precise reason for the different target of the
relative theorem.
\end{proof}

\begin{result}{F9.3}{Refuted}{Fullness into all topoi, and fullness over an algebraically closed field}
\textup{[F9; \S2, manuscript p.\ 25; scan 52.]}
Without a restriction on geometric morphisms, the proposed fullness
fails already over $\mathbb Q$. With pinned morphisms retained, it
fails over $\mathbb C$, even for smooth proper genus-one curves and
even after universal homeomorphisms are inverted.
\end{result}

\begin{proof}
For the first assertion take the $t$-adic henselian valuation of
$\mathbb Q(t)$. The coefficient action, fixing compatible roots of
$t$, gives a continuous section $G_{\mathbb Q}\to G_{\mathbb Q(t)}$
through its decomposition group. One can construct it in the Puiseux
field $\bigcup_n\overline{\mathbb Q}((t^{1/n}))$: coefficient
automorphisms preserve the algebraic closure of $\mathbb Q(t)$
inside this field and fix $t$. This section gives a geometric morphism
\[
 \mathbb Q_{\mathrm{et}}\longrightarrow
 \mathbb Q(t)_{\mathrm{et}}\longrightarrow
 (\mathbb P^1_{\mathbb Q})_{\mathrm{et}}
\]
over $\mathbb Q_{\mathrm{et}}$, whose point maps to the generic point.
An actual $\mathbb Q$-morphism $\operatorname{Spec}\mathbb Q\to
\mathbb P^1$ has closed-point image. Universal-homeomorphism
localization does not change this Hom-set, since source and target
are seminormal. This is the concrete counterexample of
Carlson--Haine--Wolf, Example 2.30, with its section constructed above.

For the second assertion choose $\lambda\in\mathbb C$ transcendental
over $\mathbb Q$ and an automorphism $\sigma$ of $\mathbb C$ with
$\sigma(\lambda)=2\lambda$. Such an automorphism is obtained by
replacing $\lambda$ by $2\lambda$ in a transcendence basis and
extending to an algebraic closure. Set
\[
 E_\lambda:y^2=x(x-1)(x-\lambda),\qquad
 E_{2\lambda}:y^2=x(x-1)(x-2\lambda).
\]
Applying $\sigma$ to the coefficients gives an isomorphism of the
underlying schemes over $\operatorname{Spec}(\sigma)$, hence an
equivalence of their \'etale topoi. Since
$\operatorname{Spec}\mathbb C_{\mathrm{et}}=\operatorname{Set}$,
this is a pinned equivalence over the specified base topos. A
$\mathbb C$-scheme isomorphism between these curves would force equality
of their $j$-invariants
\[
 j(\lambda)=\frac{256(1-\lambda+\lambda^2)^3}
 {\lambda^2(1-\lambda)^2},\qquad j(2\lambda).
\]
These rational functions differ, as the polynomial calculation in
\texttt{code/CF\_checks.py} verifies. Transcendence of $\lambda$
therefore excludes equality. A morphism inducing this equivalence
would induce a bijection of underlying topological spaces. It would
therefore be nonconstant, hence finite between these smooth proper
curves. In characteristic zero its generic degree equals the number
of geometric points of a general fibre, so that bijectivity forces
degree one. A finite degree-one map onto a normal curve is an
isomorphism, which the displayed $j$-invariants exclude.
Both curves are normal, so the
characteristic-zero universal-homeomorphism localization cannot
produce an additional isomorphism. This constructs the comparison
map between the two topoi and proves explicitly why it does not come
from a morphism of the required schemes. It does not contradict the
injectivity in F9.1.
\end{proof}

\begin{result}{F9.4}{Refuted}{A finite-field obstruction to pinned fullness}
\textup{[F9 and the characteristic question in its margin; \S2,
manuscript p.\ 25; scan 52.]}
Pinned fullness after universal-homeomorphism localization fails for
$\mathbb A^1_{\mathbb F_3}$ over $\mathbb F_3$. Thus the positive
transcendence-degree hypothesis in the positive-characteristic case
of F9.2 has an actual obstruction.
\end{result}

\begin{proof}
Put $X=\mathbb A^1_{\mathbb F_3}$ and
$p:X_{\mathrm{et}}\to B\widehat{\mathbb Z}$ for its structure map.
For any $\gamma\in\widehat{\mathbb Z}$ the central action on
$\widehat{\mathbb Z}$-sets defines a natural automorphism of the
identity of $B\widehat{\mathbb Z}$, and hence an automorphism
$\alpha_\gamma$ of $p$. The pair
$(\operatorname{id}_{X_{\mathrm{et}}},\alpha_\gamma)$ is a pinned
self-equivalence in the relative $2$-category. Choose
$\gamma\notin\mathbb Z$; for example use $2$-adic coordinate $0$
and every odd-adic coordinate $1$.

Suppose it came from a morphism in the localization. This is a
morphism $X^{\mathrm{perf}}\to X$, hence is given by
$f\in\mathbb F_3[x^{1/3^\infty}]$. Since the underlying topological
map is the identity, $f$ fixes every closed point as a point. Choose
$r$ with $P=f^{3^r}\in\mathbb F_3[x]$, of degree $D$. For each
$a\in\mathbb F_{3^N}$, preservation of its closed point gives
$P(a)=a^{3^i}$ for some $0\leq i<N$. Unless one polynomial
$P(T)-T^{3^i}$ is identically zero, their union of roots has size at
most
\[
 \sum_{i=0}^{N-1}\max(D,3^i)
 \leq ND+\frac{3^N-1}{2}<3^N
\]
for sufficiently large $N$. This contradicts the preceding assertion
for all $3^N$ elements. Therefore $f=x^{3^m}$ for an integer $m$,
where negative exponents mean inverse Frobenius on the perfection.

We check the entire $2$-isomorphism ambiguity. A natural automorphism
of $\operatorname{id}_{X_{\mathrm{et}}}$ restricts to each closed
subtopos $\operatorname{Spec}\mathbb F_{3^n,\mathrm{et}}$ because
the direct image of a closed immersion is fully faithful. Naturality
with respect to the adjunction unit $A\to i_*i^*A$ makes its
base component agree with the restriction of the original component. On that
subtopos its induced base automorphism lies in
$G_{\mathbb F_{3^n}}=n\widehat{\mathbb Z}$: the automorphisms of
the identity of $BH$ are $Z(H)$, here $H=n\widehat{\mathbb Z}$.
There is a closed point of every degree $n$ on $\mathbb A^1_{\mathbb F_3}$.
Thus its base component belongs to
$\bigcap_n n\widehat{\mathbb Z}=0$. All isomorphisms
$(\operatorname{Frob}^m)_{\mathrm{et}}\simeq\operatorname{id}$
consequently have the same base component. The canonical one, from
relative Frobenius of \'etale covers, has component $m$; with the
relative-pair convention above this represents
$(\operatorname{id},\alpha_{-m})$. It cannot represent our
$\gamma\notin\mathbb Z$. This proves the counterexample, including
the possible noncanonical natural transformations.
\end{proof}

\begin{result}{F9.5}{Known}{The proposed function-field automorphism test}
\textup{[F9; \S2, manuscript p.\ 25; scan 52.]}
Let $k$ be a number field, $C/k$ a smooth proper geometrically connected
curve of any genus, $L=k(C)$, and
$\Delta=\ker(G_L\to G_k)$. Then
\[
 \operatorname{Aut}_k(C)\longrightarrow
 \operatorname{Isom}_{G_k}(G_L,G_L)/\operatorname{Inn}(\Delta)
 \tag{F.3}
\]
is bijective, with the usual contravariance on function fields.
In particular the additional requirement that the automorphism
preserve all loop subgroups is unnecessary for this conclusion;
the automorphism obtained from (F.3) preserves them by functoriality
of valuations.
\end{result}

\begin{proof}
We use Tsujimura, Corollary 4.10, in the opened RIMS-1957 version
\cite{CF_Tsujimura2023}. It applies to function fields in one variable
over a subfield $K$ of a finitely generated extension of the field
obtained from a number field $F$ by adjoining all roots of a fixed
prime number. Take $F=K=k$; it is a subfield of that field, so every
hypothesis is met. The corollary gives the bijection between
$k$-field isomorphisms and the right-hand side of (F.3), with
conjugation exactly by the geometric kernel $\Delta$.

A $k$-automorphism of $L$ gives a rational map $C\dashrightarrow C$.
Every local ring of the smooth proper curve at a closed point is a
DVR. The valuative criterion for properness extends the rational map
there; finite presentation and separatedness glue these extensions
to a morphism. The same argument for the inverse field automorphism
gives its inverse morphism, proving the claimed bijection for $C$.
It carries the local DVR at $x$ to that at its image, and hence
carries the corresponding inertia and decomposition groups to the
groups at that image, modulo the stipulated conjugation.

On the image, scan 52 appears to abbreviate ``structure \`a lacets''
as ``\`a lcts'', rather than the transcription's alternatives
``Qcts''/``Dcts''. We record this as evidence, not a settled reading.
The proved statement covers both the version with peripheral
structure required and the stronger version where it is recovered.
\end{proof}

\subsection{F17: good reduction and the quotients of monodromy}

\begin{result}{F17.1}{Known}{The nonabelian criterion and the valid forward implications}
\textup{[F17; \S3, manuscript pp.\ 49--51; scans 100--103.]}
Let $R$ be a DVR with fraction field $K$ and residue characteristic
$p\geq0$. Let $U=\bar U\setminus D$ be a smooth geometrically connected
hyperbolic curve over $K$, with $\bar U$ smooth proper and $D$ reduced
and geometrically \'etale. Good reduction means a smooth proper model
of $\bar U$ over $R$ on which the boundary extends to a finite
\'etale divisor. Good reduction is equivalent to trivial inertia in
$\operatorname{Out}(\Pi^{(p')})$; it is also equivalent to trivial
inertia in $\operatorname{Out}(\Pi^{(\ell)})$ for one, equivalently
every, prime $\ell\ne p$. For $p=0$ the pro-prime-to-$p$ group is the
full geometric group, so this proves exactly a)$\Leftrightarrow$b).

Good reduction implies c) and d) with prime-to-$p$ coefficients, through
the characteristic quotient maps
$\Pi^{(p')}\to\Pi^{(p'),\mathrm{ab}}\to H_1(U_{\bar K},\mathbb Z_\ell)$.
When the residue field is finitely generated over $\mathbb Q$, a) and
b) imply e) for every splitting of the residue exact sequence: the
resulting outer action is that of the actual smooth special fibre.
\end{result}

\begin{proof}
The precise source is Tamagawa, Theorem (5.3), printed p.\ 177,
equivalently Theorem (0.8), with the definitions immediately before
it \cite{CF_Tamagawa1997}. Its hypotheses are a DVR, a smooth proper
geometrically connected curve and relatively \'etale effective boundary
over its fraction field, and $2-2g-r<0$. It makes the four conditions
just stated equivalent; no semistability assumption is added.
Thus its proper-curve case includes Oda's criterion and its marked
case covers colliding cusps as well.

Inner automorphisms act trivially on abelianization. Therefore an
outer inertia action equal to the identity induces the identity on
each indicated abelian quotient, giving the actual comparison maps
for a)$\Rightarrow$b)$\Rightarrow$c)$\Rightarrow$d).
In good reduction, specialization identifies the geometric groups
of generic and special fibres, and identifies their arithmetic
extensions after pullback along $G_K\to G_k$. The outer action
therefore factors through $G_k$. Composing this factorization with
any continuous splitting $G_k\to G_K$ yields that same special-fibre
action. Since $\Pi$ is centre-free, C4 identifies the extension
reconstructed from the action with the extension of the special
fibre. This proves e), without an additional field extension.

In mixed characteristic one cannot include the residue prime in the
homological conclusion. For example the smooth marked curve
$\mathbb P^1_{\mathbb Z_p}\setminus\{0,1,\infty\}$ has good
reduction, but its generic $H_1$ contains $\mathbb Z_p(1)^2$ and
$p$-adic cyclotomic inertia is nontrivial. This is why the criterion
above specifies its quotient rather than silently using the full
group at residue characteristic $p$.
\end{proof}

\begin{result}{F17.2}{Refuted}{A compact-type genus-two counterexample, including class two}
\textup{[F17; \S3, manuscript pp.\ 49--51; scans 100, 102--103.]}
Over $R=\mathbb Q[[t]]$, the smooth projective completion of
\[
 C_t:\quad y^2=(x^3-t^6)(x^3-1)                         \tag{F.4}
\]
has genus two, bad reduction, and trivial inertia on all of its
abelianizations. Its inertia action is even trivial on every
class-two nilpotent geometric quotient. Its full outer inertia
action is nontrivial. Thus c)$\Rightarrow$a) and d)$\Rightarrow$a)
are false for proper curves.
\end{result}

\begin{proof}
The sextic discriminant is
$729t^{12}(1-t^6)^6$, as checked in the attached program. It is
nonzero in $\mathbb Q((t))$, so the double cover has six distinct
branch points and genus $(6-2)/2=2$.

We describe its stable model explicitly. Separate the two sets of
three branch points using the model of the base projective line with
local plumbing equation $xz=t^2$. On the inner chart put
$x=t^2X,y=t^3Y$; the equation becomes
\[
 Y^2=(X^3-1)(t^6X^3-1),
\]
with special fibre $E_{\mathrm{in}}:Y^2=1-X^3$. On the outer
component the normalization is
$E_{\mathrm{out}}:V^2=x(x^3-1)$, with $V=y/x$.
Both are smooth genus-one curves.
On the outer curve the substitutions $X=1/x$ and $Y=V/x^2$
give $Y^2=1-X^3$, so the two smooth proper components are
isomorphic over $\mathbb Q$. Over $\mathbb Q(\zeta_3)$ the
automorphism $(X,Y)\mapsto(\zeta_3X,Y)$ fixes the origin and
has order three, giving the CM action by $\mathbb Z[\zeta_3]$.
At the joining point set
$A=(1-z^3)(x^3-1)$, a unit. The equations are
\[
 V^2=xA,\qquad W^2=z/A,\qquad VW=t,
 \qquad W=t/V.
\]
They follow by substitution in (F.4); conversely they recover the
original equation on the generic fibre. Since $A(0,0)=-1$, the
implicit-function theorem for complete local rings identifies the
completed local equation with $VW=t$. The model is smooth away
from this node, since the remaining branch sections are disjoint
and simple. Normalizing the proper separated-branch base model
therefore gives a proper stable curve whose special fibre consists
of these two elliptic curves meeting transversely at one point.
Both components are stable, and the fibre is not smooth. Uniqueness
of the stable model (the input also used in Tamagawa, Remark (5.2))
excludes a smooth model.

The single vanishing circle separates the two genus-one pieces.
Choose a geometric topological presentation
\[
 \pi_1(C_t(\mathbb C))=
 \langle a_1,b_1,a_2,b_2\mid[a_1,b_1][a_2,b_2]=1\rangle,
 \qquad c=[a_1,b_1].
\]
The plumbing $VW=t$ twists one boundary circle once as $t$ circles
zero. Cutting at this circle shows directly that the induced
automorphism conjugates $a_1,b_1$ by $c$ and fixes $a_2,b_2$, or
its inverse with the opposite loop orientation. This description
can also be obtained by identifying $V$-annulus and $W$-annulus
coordinates through $W=t/V$: a full turn of $t$ adds one turn to
the gluing coordinate. Finite topological covers and finite \'etale
covers are identified by SGA~1 XII, Theorem 5.1; the same formula
therefore describes the profinite and pro-$\ell$ tame inertia actions.
The family is defined algebraically near $t=0$, so the topological
loop maps to the positive generator of its tame inertia, with this
identification.

The element $c$ belongs to the second lower-central subgroup
$\Gamma_2$. In the class-two quotient $\Pi/\Gamma_3$ it is central,
so the displayed automorphism is the identity there, and in
particular on $\Pi^{\mathrm{ab}}$. This is an identity in the
quotient itself, not merely an outer identity. The full outer
action is nontrivial by F17.1 and the proved failure of good
reduction. The special fibre is of compact type and its Jacobian
is the product of the two elliptic component Jacobians. More
precisely, the relative $\operatorname{Pic}^0$ of a compact-type
model is an abelian scheme, and its special fibre is the product
of the component Jacobians; the construction and properness
argument are given in \cite[\S1, p.\ 1]{CF_DupuyRabinoff2017}.
That text assumes algebraically closed residue field. Apply it
after the faithfully flat strict-henselian base extension;
formation of the relative Picard functor commutes with that base
change, and smoothness and properness descend. This proves the
same assertion over $\mathbb Q[[t]]$ and explains the absence
of a homological vanishing cycle.
\end{proof}

\begin{result}{F17.3}{Refuted}{Colliding cusps detected by class-two monodromy}
\textup{[F17; \S3, manuscript pp.\ 50--51; scans 102--103.]}
Let
\[
 U_t=\mathbb P^1_{\mathbb Q((t))}\setminus\{0,t,1,\infty\}.
 \tag{F.5}
\]
It has bad marked reduction and trivial inertia on
$\Pi^{\mathrm{ab}}=\widehat{\mathbb Z}(1)^3$, while inertia acts
nontrivially in the outer automorphism group of its class-two
pro-$\ell$ quotient for every prime $\ell$.
\end{result}

\begin{proof}
All four punctures are rational. The homology is the quotient of
the free Tate module on the four punctures by its diagonal. The
$t$-adic inertia fixes all roots of unity (they are constants), so
its action on this quotient is trivial. The cross ratio of the
ordered four points is $t$ up to one of its six cross-ratio
transforms. In a smooth marked projective line over a strictly
henselian trait, carrying three disjoint sections to $0,1,\infty$
makes the fourth coordinate and its difference from $1$ units.
Our valuation $v(t)>0$ precludes this. Alternatively, blowing up
the collision at $(x,t)=(0,0)$ gives two three-pointed projective
lines joined at a node, the stable model. This proves failure of
good reduction independently of the monodromy computation.

Choose free geometric generators $a,b,c$ around $0,t,1$; the loop
at infinity is $(abc)^{-1}$. A loop of $t$ around zero is the full
twist of the first two punctures. Cutting out a disc containing
exactly these two points shows its action to be
\[
 a\longmapsto(ab)a(ab)^{-1},\qquad
 b\longmapsto(ab)b(ab)^{-1},\qquad c\longmapsto c,             \tag{F.6}
\]
up to reversing the generator of inertia. This is the same local
plumbing calculation as in F17.2, now for a disc with two punctures.
Let $[x,y]=xyx^{-1}y^{-1}$ and pass to the free class-two pro-$\ell$
group. Its commutator subgroup is freely generated as a
$\mathbb Z_\ell$-module by
$c_{ab}=[a,b],c_{ac}=[a,c],c_{bc}=[b,c]$. Formula (F.6) becomes
\[
 a\mapsto a c_{ab}^{-1},\qquad
 b\mapsto b c_{ab},\qquad c\mapsto c.                       \tag{F.7}
\]
Conjugation by an element with abelian coordinates $(u,v,w)$ gives
\[
 a\mapsto a c_{ab}^{-v}c_{ac}^{-w},\quad
 b\mapsto b c_{ab}^{u}c_{bc}^{-w},\quad
 c\mapsto c c_{ac}^{u}c_{bc}^{v}.
\]
Matching the first two images forces $u=v=1,w=0$; the third then
has the nonzero commutator coordinates $(1,1)$. Thus (F.7) is not
inner. This calculation holds over every $\mathbb Z_\ell$, including
$\ell=2$. The program verifies it separately in the finite class-two
exponent-five quotient by examining all $5^3$ possible abelian
coordinates of an inner automorphism. This finite check supplements
the integral proof and does not replace its all-primes conclusion.
\end{proof}

\begin{result}{F17.4}{Refuted}{Homological unramifiedness need not give a geometric residue action}
\textup{[F17; \S3, manuscript pp.\ 50--51; scans 101--103, condition e).]}
For the curve (F.5), choose the coefficient splitting
$G_{\mathbb Q}\to G_{\mathbb Q((t))}$ which fixes compatible roots
of $t$. The resulting outer action on $\widehat F_3$ is not realized
by a smooth hyperbolic curve over any finite extension of
$\mathbb Q$. Thus neither c) nor d) implies e), even allowing the
finite extension in e).
\end{result}

\begin{proof}
We use the nonlinear inertia criterion of Nakamura,
Theorem (2.1.1) \cite{CF_Nakamura1994}. Its hypotheses are a smooth
geometrically connected affine hyperbolic curve over a number field
and a full class of finite groups; we take the class of all finite
groups. It states that a nonidentity geometric element belongs to
an inertia group precisely when it has a closed cyclotomic normalizer
whose image in the arithmetic Galois group is open. ``Cyclotomic''
here means that the action on its closed procyclic subgroup factors
through that arithmetic image and is the cyclotomic character.

Let $E_0$ be the extension of $G_{\mathbb Q}$ by $\Pi=\widehat F_3$
obtained by the chosen coefficient splitting. Each of the four
cusps supplies a decomposition subgroup in $E_0$ and a tangential
section normalizing its loop through the cyclotomic character.
There is a fifth such subgroup for the vanishing loop $d=ab$.
We justify this extra subgroup locally. At the node of the stable
marked model the regular equation is $xz=t$. The punctured formal
neighborhood, with the two vertical divisors removed, has tame
geometric group $T^2$, by the two-parameter form of SGA~1 XIII,
Proposition 5.2. Its map to the tame inertia of the base is
\[
 T^2\longrightarrow T,\qquad (r,s)\longmapsto r+s,
\]
because $t=xz$. Coefficient automorphisms split its arithmetic
extension and act cyclotomically on both factors. Choose the roots
so that $x^{1/n}z^{1/n}=t^{1/n}$ for every $n$; coefficient
automorphisms fix both factors and hence their product. They
therefore map to our specified coefficient splitting of the base. The kernel
$\{(r,-r)\}$ maps to the vanishing loop subgroup $\langle ab\rangle$
of the generic geometric curve. The map on this kernel is injective:
$ab$ has primitive abelianization in the free geometric group, and
the local annulus identifies the two orientations as opposite.
Consequently its image together with the coefficient section is a
closed subgroup of $E_0$ projecting onto $G_{\mathbb Q}$ and acting
on $\langle ab\rangle$ cyclotomically. This construction persists
after restriction to every finite number field.

Suppose $E_0$ became the arithmetic extension of a smooth hyperbolic
curve $V/k$ after such a restriction. Its geometric group has
abelian rank three, so $V$ is affine and
$2g(V)+r(V)-1=3$, whence $r(V)\leq4$. Transfer the five cyclotomic
normalizers through the asserted extension isomorphism. Nakamura's
criterion puts the five elements
\[
 a,\ b,\ c,\ (abc)^{-1},\ ab
\]
in inertia groups of $V$. In $H_1(\Pi,\mathbb Q_\ell)$ their lines
are generated by
\[
 e_1,\quad e_2,\quad e_3,\quad-e_1-e_2-e_3,\quad e_1+e_2.
\]
They are five distinct lines. Elements belonging to conjugate
inertia groups at one cusp have proportional abelian classes.
Hence these five elements require at least five distinct cusps of
$V$, contradicting $r(V)\leq4$. The attached program checks the
pairwise rank-two test for these five lines. This proves the failure
of e) with its full finite-extension allowance.
\end{proof}

\begin{result}{F17.5}{Refuted}{A geometric residue action can coexist with bad reduction}
\textup{[F17 and an additional flagged step; \S3, manuscript pp.\ 50--51;
scans 101, 103.]}
Over $K=\mathbb Q((t))$ let
\[
 E_t:y^2=x^3-t^2x,\qquad U_t=E_t\setminus\{O\}.
 \tag{F.8}
\]
It has bad reduction. For every continuous splitting
$s:G_{\mathbb Q}\to G_K$, however, its restricted outer action is
realized by a once-punctured elliptic curve over $\mathbb Q$ itself.
Thus e)$\Rightarrow$a) is false even with ``every splitting'', and
bad reduction alone does not force the proposed residue action to
be nongeometric.
\end{result}

\begin{proof}
Write $E_0:Y^2=X^3-X$. Over $K(\sqrt t)$ the substitution
$x=tX$, $y=t\sqrt t\,Y$ gives an isomorphism $E_0\simeq E_t$,
as checked by direct substitution. The nonidentity automorphism of
this quadratic extension acts through $[-1]$ on $E_0$. Therefore
$E_t$ is the twist by the ramified quadratic character
$\epsilon_t:G_K\to\{\pm1\}$ associated to $t$. Its inertia action
on $H_1(E_{\bar K},\mathbb Z_\ell)$ is $-1$, so good reduction
is impossible by the forward implication in F17.1 (or smooth
proper base change).

For any splitting $s$, the composite $\epsilon_t\circ s$ is a
quadratic character of $G_{\mathbb Q}$, hence is the Kummer
character of some $d\in\mathbb Q^*/\mathbb Q^{*2}$. Restriction
of the twisting cocycle is precisely the cocycle defining the
quadratic twist $E_0^{(d)}$. Functorial descent of finite \'etale
covers then identifies the pulled-back extension with that of
$E_0^{(d)}\setminus\{O\}$. Explicitly, the geometric action is
$g\mapsto\iota(\epsilon_t(s(g)))\rho_0(g)$, where
$\iota(+1)=\operatorname{id}$ and $\iota(-1)=[-1]$, and this is the descent action for that twist.
This proves e) without enlarging the residue field. The point
$(0,0)$ is fixed by $[-1]$, so the same construction can be made
with actual point sections, rather than only outer actions.
\end{proof}

\begin{result}{F17.6}{Proved}{Stable bad reduction obstructs every geometric residue action}
\textup{[F17; \S3, manuscript pp.\ 49--51; scans 100--103;
supplied R4 and H2.]}
Let $R$ be a henselian discrete valuation ring with fraction field
$K$ and residue field $k$ finitely generated over $\mathbb Q$.
Let $U/K$ be a hyperbolic curve whose smooth proper compactification
and reduced boundary have a stable marked model $(\mathcal C,\mathcal D)$
over $R$. Put $\Pi=\pi_1(U_{\bar K})$, with its full profinite topology.
For a continuous section $s:G_k\to G_K$ of the residue map, let
\[
 E_s=\pi_1(U)\mathbin{\times}_{G_K,s}G_k.
\]
Every geometric node of $\mathcal C_k$ gives an injected subgroup
$I_e\simeq\widehat{\mathbb Z}(1)$ of $\Pi$, well defined up to
$\Pi$-conjugacy. In $E_s$ it is essentially cyclotomically normalized:
after a finite extension of $k$ fixing that node and its two branches,
a closed subgroup of its normalizer maps onto the residue Galois
group and acts on $I_e$ by the cyclotomic character. This holds for
every $s$. Distinct geometric nodes and cusps give distinct conjugacy
classes of maximal procyclic subgroups.

If $\mathcal C_k$ has a node, the extension $E_s$ is not isomorphic,
over $G_k$, to the arithmetic fundamental-group extension of any
smooth hyperbolic curve over $k$, and it does not become so after
any finite extension of $k$. If $r>0$ is the number of geometric
cusps, the number of conjugacy classes of essentially cyclotomically
normalized maximal procyclic subgroups of $\Pi$ equals $r$ exactly
when this stable marked model is smooth. In the singular case it is
at least $r+n$, where $n>0$ is the number of geometric nodes.
\end{result}

\begin{proof}
We give the full comparison which was missing in the supplied prior
return E4. All logarithmic fundamental groups below classify finite
Kummer \'etale covers; all sets of primes are the set of \emph{all}
primes. This choice is legitimate because $k$ has characteristic
zero.

First pass from $R$ to its completion. The restriction map
$G_{\widehat K}\to G_K$ is compatible with the residue maps to
$G_k$. The two inertia groups are $T=\widehat{\mathbb Z}(1)$:
over the maximal unramified extension every finite extension is
killed by adjoining a root of a uniformizer, by the tame local
calculation of SGA~1 XIII, Propositions 5.1--5.2. The restriction
map induces the identity on $T$, as is seen on the same compatible
roots of that uniformizer, and the identity on $G_k$. It is therefore
an isomorphism of the two extensions. Geometric fundamental groups
of a smooth curve are invariant under an extension of algebraically
closed fields in characteristic zero: this follows from SGA~1 XIII,
Proposition 4.3, Examples 4.4, and \S4.5.3, applied to the
constant smooth proper pair with a chosen geometric point section.
The base-change map
$\pi_1(U_{\widehat K})\to\pi_1(U)$ consequently induces
isomorphisms on both geometric kernel and arithmetic quotient,
and hence is an isomorphism. Thus it transports each specified
$s$, its pullback extension, and the stable special fibre without
changing the question.

We now work over the complete ring. Give $S=\operatorname{Spec}R$
its divisorial log structure, and $\mathcal C$ the log structure
from $\mathcal C_k\cup\mathcal D$. These are fs log schemes;
$\mathcal C^{\log}$ is log regular and proper over $S^{\log}$.
For example a split node of thickness $m$ has completed equation
$xy=t^m$ after absorbing a unit into $y$, and monoid
\[
 P=\langle X,Y,T_0\mid X+Y=mT_0\rangle.
\]
Its group completion has rank two; the closed log stratum is
regular of dimension zero, and the local ring has dimension two.
These are precisely the log-regularity dimension and regularity
conditions. At a smooth point and at a marked section the charts
are respectively the one-parameter and two-parameter normal-crossing
charts. The monoid $P$ is saturated: in the basis $(X,T_0)$ its
integral cone is generated as a monoid by $(1,0),(0,1),(-1,m)$.

We use the following two opened, numbered comparison results.
Illusie, Theorem 7.6 \cite{CF_Illusie2002}, states that for a
log-regular fs scheme, restriction to its trivial-log open
identifies finite Kummer \'etale covers with finite \'etale covers
tame along the boundary at its codimension-one points.
Here every such residue characteristic is zero, so all covers are
tame. Hoshi, Corollary 3.9 \cite{CF_Hoshi2009}, states that a
proper log-regular fs scheme over the spectrum of a complete local
ring has the same finite Kummer \'etale cover category as its
strict closed log fibre. Applying these results also to the base
gives a commuting square of isomorphisms
\[
 \begin{array}{ccc}
 \pi_1(\mathcal C_k^{\log})&\xrightarrow{\sim}&\pi_1(U)\\
 \downarrow&&\downarrow\\
 \pi_1(k^{\log})&\xrightarrow{\sim}&G_K .
 \end{array}                                                    \tag{F.9}
\]
The upper arrow is the strict-closed-fibre comparison followed by
the inverse of restriction to the interior. These functors commute
with finite extension of the residue field, with log base change,
and with the morphism from a strict point at a node. The resulting
maps of fundamental groups have only the usual base-point path,
hence inner-conjugacy, ambiguity.

For clarity, the geometric kernel in (F.9) is also identified in
an opened source. Mochizuki, Example 2.5
\cite{CF_MochizukiComb2007}, treats a stable log curve over the
standard log point $\mathbb N\to k$, $1\mapsto0$. Its log
Galois group is an extension of $G_k$ by $T$, and the geometric
kernel is the fundamental group $\Pi_{\mathcal G}$ of the
semi-graph of groups of the stable curve. The identification is
natural and the arithmetic action preserves that graph.
The finite covers in this description are obtained by matching
the tame branch covers at each node, with inverse branch
orientations; the chart $X+Y=mT_0$ gives exactly this matching.
Hoshi--Mochizuki, Definition 2.10(iv), Lemma 2.11(ii)--(iii), and
Remark 2.14.1 \cite{CF_HoshiMochizukiNodes}, make the log-group
injection and its change under the smoothing-exponent map
explicit. In particular, with all primes, (F.9) identifies
\[
 \Pi_{\mathcal G}
 =\ker[\pi_1(\mathcal C_k^{\log})\to\pi_1(k^{\log})]
 \xrightarrow{\sim}\ker[\pi_1(U)\to G_K]=\Pi .                 \tag{F.10}
\]
Thus (F.10) is an isomorphism of the full geometric groups, with
the arithmetic action inherited from (F.9); it is not only an
isomorphism of their abelianizations or pro-$\ell$ quotients.

Fix a geometric node. Enlarge $k$ by a finite extension which
defines the node and its two branches, and make the corresponding
unramified extension of the trait. This replaces $G_k$ by its
specified open subgroup and $s$ by its restriction. The log point
at that node has arithmetic group
\[
 D_e=T_P\rtimes G_k,\qquad
 T_P=\{(a,b,c)\in T^3:a+b=mc\},
\]
where $G_k$ acts cyclotomically on all three coordinates. Its map
to the base log Galois group, using the same uniformizer splitting,
is
\[
 (a,b,c,g)\longmapsto(c,g)\in T\rtimes G_k=G_K.                \tag{F.11}
\]
This map is surjective, since $(mc,0,c)\in T_P$ for every $c$.
Its kernel is $\{(a,-a,0)\}\simeq T$. Compatible Kummer roots
may be chosen with $x^{1/n}y^{1/n}=(t^{1/n})^m$, which proves
(F.11) and its compatibility with the specified base action.
The map from this kernel into (F.10) is the injected edge group of
Lemma 2.11. Stability is essential here: these are the edges of
the stable curve, not contractible circles attached to an unstable
rational tail. Resolving a node of thickness $m$ merely subdivides
its edge and gives the same essential circle.

Let $H_s$ be the inverse image of $s(G_k)$ under (F.11).
It is closed, maps onto $s(G_k)$, and its geometric kernel is the
displayed antidiagonal $T$. Its image in $E_s$ therefore is a
closed normalizing subgroup mapping onto $G_k$. Conjugation on
that kernel is $a\mapsto\chi(g)a$, since $T_P$ is abelian and
its arithmetic action is cyclotomic. No complement in $H_s$ has
been chosen or assumed. This proves the assertion for every
splitting, including one different from the uniformizer splitting
used to write (F.11).

Mochizuki, Proposition 1.2(i)--(ii)
\cite{CF_MochizukiComb2007}, applies to every semi-graph of
profinite pointed-stable-curve type: its edge subgroups are
commensurably terminal, and open intersection of two edge
subgroups forces the same edge. Thus each edge group is maximal
procyclic, since an abelian group containing it normalizes it and
is contained in its commensurator. A conjugacy between groups
belonging to two geometric edges would identify those edges by
part (i); distinct nodes and cusps therefore give distinct
$\Pi$-conjugacy classes. Marked sections give the same
cyclotomic-normalizer construction by their one-branch Kummer
charts. This proves the bound $r+n$. If $n=0$, (F.9) is the
good-reduction comparison, and Stix, Theorem 20 and Corollary 21,
give exactly the $r$ cusp classes. The required hypothesis that
$k$ distinguishes tori from abelian varieties follows from the
checked finitely generated characteristic-zero result AE1 of
\cite{PriorCyclotomic}: a Tate-module map from a split torus
would have image fixed by the cyclotomic kernel. This also proves
the asserted counting criterion.

We finally exclude realization by an arbitrary smooth curve,
whose number of cusps has not been fixed in advance. This needs
an additional argument in the affine case. The comparison (F.10)
gives the graphic quotient and the homology filtration described
in Mochizuki, Definition 1.1(ii) and Remarks 1.1.3--1.1.4.
If the dual graph has first Betti number $b>0$, killing its vertex
groups gives a surjection
\[
 \Pi\twoheadrightarrow\widehat{\pi_1(\mathcal G)},\qquad
 H_1(\Pi,\mathbb Q_\ell)\twoheadrightarrow\mathbb Q_\ell^b.
                                                               \tag{F.12}
\]
After a finite extension fixing the graph, the right-hand module
has trivial $G_k$-action. This statement is independent of $s$:
logarithmic inertia fixes the underlying graph, and the entire
action on this quotient factors through the finite action of
$G_k$ on that graph.

No smooth curve $V/k$ has a nonzero trivial quotient of its
$\mathbb Q_\ell$-homology, after any finite extension of $k$.
Indeed, after its cusps split, the actual homology sequence is
\[
 0\longrightarrow \mathbb Q_\ell(1)^{r_V-1}
 \longrightarrow H_1(V_{\bar k},\mathbb Q_\ell)
 \longrightarrow V_\ell\operatorname{Jac}(\bar V)
 \longrightarrow0                                             \tag{F.13}
\]
for $r_V>0$; in the proper case the first term is zero.
A map from the first term to the trivial representation is zero,
because the $\ell$-adic cyclotomic character has infinite image
over every field finitely generated over $\mathbb Q$.
For the last term put $A=\operatorname{Jac}(\bar V)$. The Weil
pairing $T_\ell A\times T_\ell(A^\vee)\to\mathbb Z_\ell(1)$
identifies $(V_\ell A)^*$ with $V_\ell(A^\vee)(-1)$. Its invariants
under the cyclotomic kernel vanish by the same AE1 theorem,
applied to the dual abelian variety. Hence its $G_k$-invariants
vanish as well. Applying $\operatorname{Hom}_{G_k}(-,\mathbb Q_\ell)$
to (F.13) proves the assertion and contradicts (F.12).

It remains to treat a tree. For a proper curve, the existence of
one node subgroup already contradicts Stix's Theorem 20 for a
proper realizing curve. Such a realizing curve must indeed be
proper: the continuous group $H^2(\Pi,\mathbb Z/\ell)$ is
$\mathbb Z/\ell$ for a proper hyperbolic curve and zero for an
affine one, by the curve $K(\pi,1)$ comparison and the surface/free
group presentations. An isomorphism of geometric groups induces
an isomorphism on this cohomology.

For an affine curve with $r>0$ and tree dual graph, the same graph
presentation gives an exact sequence of $G_k$-modules
\[
 0\longrightarrow\mathbb Q_\ell(1)^{r-1}
 \longrightarrow H_1(\Pi,\mathbb Q_\ell)
 \longrightarrow\bigoplus_v V_\ell\operatorname{Jac}(C_v)
 \longrightarrow0,                                            \tag{F.14}
\]
after a finite extension defining the smooth normalized components
$C_v$, nodes, branches, and marks. Here is the rank computation:
at each vertex the sum of its boundary classes is zero; eliminating
an end vertex of the tree expresses its unique incident node
class as a sum of its cusp classes. Successively removing end
vertices expresses every node class in terms of the $r$ cusp
classes, with only their single total-sum relation. Quotienting
by these classes fills every puncture on every component and
gives precisely the last term of (F.14). The action there is the
ordinary arithmetic action of the actual curves $C_v/k$; log
inertia is trivial on this quotient. These facts follow from the
same functorial vertex maps in (F.10), so they hold after every
restriction along $s$.

The torus/abelian-variety vanishing used above gives
\[
 \dim_{\mathbb Q_\ell}
 \operatorname{Hom}_{G_k}\bigl(\mathbb Q_\ell(1),
 H_1(\Pi,\mathbb Q_\ell)\bigr)=r-1.                           \tag{F.15}
\]
A smooth realizing curve must be affine by $H^2$ as above,
including the distinction between $r=1$ and $r=0$.
Applying (F.13) to it shows that the same dimension is $r_V-1$;
hence $r_V=r$. Stix's criterion then gives it exactly $r$ maximal
essentially cyclotomic classes, contradicting the $r+n$ already
constructed. Every argument survives a further finite extension
of $k$. This proves the nonrealization assertion in its full
finite-extension form and, in particular, settles the two
isomorphic CM elliptic components of F17.2.
\end{proof}

\begin{result}{F17.7}{Proved}{Condition e detects potential good reduction}
\textup{[F17; \S3, manuscript pp.\ 50--51; scans 101--103;
the correction proposed in the reader's notes.]}
Let $R,K,k,U$ have the hypotheses of F17.6 except that a stable
model over $R$ is not assumed. Fix any continuous residue
splitting $s:G_k\to G_K$. The following three assertions are
equivalent: $U$ has good marked reduction after a finite extension
of $K$; the outer action obtained through $s$ is realized by a
smooth hyperbolic curve over $k$; that action is realized after
restriction to $G_{k'}$ for some finite extension $k'/k$.
Consequently ``some splitting'' and ``every splitting'' give
the same condition e), and allowing a finite residue extension
does not enlarge it in these hypotheses.
\end{result}

\begin{proof}
We first make the finite-extension comparison explicit. If
$K_1/K$ is finite inside the chosen separable closure, put
\[
 H=s^{-1}(G_{K_1})\subset G_k,
 \qquad G_L=G_{K_1}\cap p^{-1}(H),\qquad p:G_K\to G_k.
                                                               \tag{F.16}
\]
The subgroup $H$ is open, so $H=G_{k_2}$ for a finite extension
$k_2/k$. It is contained in the image of $G_{K_1}$ in $G_k$.
The fixed field $L$ of the displayed open subgroup is finite
over $K$, contains $K_1$, and has residue field $k_2$: the
restriction $p:G_L\to H$ is surjective, because every $h\in H$
has a lift in $G_{K_1}$. Moreover $s(H)\subset G_L$, and
$s|_H$ is a section of that residue map. The arithmetic
fundamental-group base-change square for a field extension gives
\[
 \pi_1(U_L)=\pi_1(U)\times_{G_K}G_L,
 \qquad
 \pi_1(U_L)\times_{G_L,s|_H}H=E_s\times_{G_k}H.                \tag{F.17}
\]
These are actual continuous extension isomorphisms, up to
geometric inner conjugation from base points.
Residue splittings exist in this characteristic-zero setting:
write a separable closure as the maximal unramified extension
with all compatible roots of a uniformizer adjoined, and extend
each residue automorphism by fixing those roots. The Kummer
description used in F17.6 proves the resulting section is
continuous.

Suppose $U_{K_1}$ has good marked reduction. We prove realization
over $k$ itself, including its descent map. Enlarge $K_1/K$ to
a finite Galois extension, which preserves good reduction, and
write $G=\operatorname{Gal}(K_1/K)$ and $k_1$ for its residue
field. Uniqueness of the stable marked model extends the natural
semilinear $G$-action on the generic pair to its smooth marked
model: for each $\sigma\in G$ the generic isomorphism
$\sigma^*(\bar U_{K_1},D_{K_1})\to(\bar U_{K_1},D_{K_1})$
extends uniquely between the two stable models. The cocycle
identities hold because they hold on the generic fibre and the
extension is unique. The marked divisor is unordered here;
permuting labels merely relabels the ordered stable model and
does not obstruct this uniqueness argument. Reducing gives isomorphisms
$\lambda_\sigma:\sigma^*(C_1,D_1)\to(C_1,D_1)$ of the
special pair over $k_1$, where $\sigma$ acts on $k_1$ by its
residue automorphism.

The continuous homomorphism
$\alpha:G_k\xrightarrow{s}G_K\to G$ has finite image and
projects to the ordinary restriction map on $k_1$.
Its kernel is $G_{k_2}$ for a finite Galois extension $k_2/k$
containing $k_1$. For $\gamma\in\operatorname{Gal}(k_2/k)$,
the base change of $\lambda_{\alpha(\gamma)}$, combined with
the scalar action of $\gamma$ on $k_2$, defines an isomorphism
\[
 \gamma^*(C_1,D_1)_{k_2}\longrightarrow(C_1,D_1)_{k_2}.
                                                               \tag{F.18}
\]
Here $\alpha(\gamma)$ means the value of $\alpha$ on any lift;
the definition is independent of that lift by the choice of
$k_2$. The identities for the $\lambda_\sigma$ make (F.18) an
actual finite Galois descent datum.

We also verify its effectiveness. The line bundle
$\omega_{C_1}(D_1)^{\otimes3}$ is very ample on every smooth
hyperbolic marked curve: its degree is at least $2g+1$ when
$g\ge1$, and positive when $g=0$. Its formation is functorial
in the pair, so (F.18) gives a compatible linear action on its
space of sections. For a finite Galois extension $L/k$ and a
finite-dimensional semilinear $L$-space $M$, the natural map
$L\otimes_k M^{\operatorname{Gal}(L/k)}\to M$ is an
isomorphism. Explicitly, trace-dual bases $a_i,b_i$ of $L/k$
give invariant vectors $v_i=\sum_\gamma\gamma(b_i v)$ and
$v=\sum_i a_i v_i$; invertibility of the embedding matrix
$(\gamma(a_i))$ gives injectivity. Apply this to the section
space and to each graded piece of the stable homogeneous ideals
of the embedded curve and its marked divisor. They descend to
a projective marked pair $(C_s,D_s)$ over $k$, whose base
change to $k_2$ is the displayed pair. Smoothness and the
finite-\'etale property of the marks descend by this faithfully
flat field extension. Put $V_s=C_s\setminus D_s$.

Here is the comparison of the full arithmetic actions, including
the action of the infinite subgroup $G_{k_2}$. Put
\[
 Q=G\times_{\operatorname{Gal}(k_1/k)}G_k,
 \qquad q:G_K\longrightarrow Q,
 \quad \delta\longmapsto(\delta|_{K_1},p(\delta)).
\]
The map $q$ is surjective: lift the first coordinate to $G_K$
and adjust its residue by an element of $G_{K_1}$, whose residue
map surjects onto $G_{k_1}$. Its kernel is
$G_{K_1}\cap I_K=I_{K_1}$. Good reduction over $K_1$ makes this
kernel act trivially in $\operatorname{Out}(\Pi)$, so $\rho_U$
factors through $Q$. On the geometric special pair, $(\sigma,g)$
acts through $\lambda_\sigma$ together with $g$ on the scalars
$\bar k$; equality of their restrictions to $k_1$ makes this
action well defined. The good-reduction comparison for covers
is functorial in isomorphisms of marked models. It therefore
identifies this $Q$-action on the special geometric group with
the factored generic outer action. Finally
$q\circ s(g)=(\alpha(g),g)$, exactly the action obtained from
the descent datum (F.18). Thus the descent action of $G_k$ on
$\pi_1((V_s)_{\bar k})$ is identified with $\rho_U\circ s$ on
$\Pi$. In particular, on $G_{k_2}=\ker\alpha$ the second
coordinate retains the ordinary arithmetic action on the curve
over $k_2$; no factorization of that action through a finite
group is being asserted.
Since $\Pi$ is centre-free, the extension reconstruction in C4
identifies their full arithmetic extensions over $G_k$ as well.
This proves realization over $k$ for the arbitrary chosen
splitting. Realization after a finite residue extension follows
by base change.

For the converse choose $K_1/K$ over which the marked curve has
stable reduction. The precise source is Knudsen, Theorem 2.7,
printed p.\ 179 \cite{CF_Knudsen1983}: for every
$2g-2+r>0$ the stack of stable $r$-pointed genus-$g$ curves is
proper and smooth over $\mathbb Z$. After first splitting and
ordering the finitely many geometric cusps, its properness,
applied to the classifying generic point, supplies the stable
marked model after a finite trait extension; separatedness gives
uniqueness and compatibility with further base change. The
unmarked semistable reduction statement is also Stacks,
Theorem 55.18.1 \cite{CF_StacksStable}.
If the geometric stable special fibre has a node, it retains that
node after every finite base change: base change of a stable
marked curve is stable, and its special fibre is the corresponding
field extension of that nodal fibre. Use (F.16)--(F.17), and
intersect $H$ with any open residue subgroup over which the
proposed realizing curve is defined. F17.6 then contradicts
that realization. Thus the stable special fibre is smooth,
which says exactly that $U_{K_1}$ has good marked reduction.
\end{proof}

\begin{result}{F17.8}{Refuted}{Individual node invariance needs the finite-extension qualification}
\textup{[F17; \S3, manuscript pp.\ 50--51; scans 101--103;
the literal individual-node assertion in supplied H2(b).]}
The full residue Galois group can permute the node groups. The
strongest uniform conclusion used in F17.6 is invariance under
an open subgroup fixing the node and its branches, with cyclotomic
action there.
\end{result}
\begin{proof}
Over $\mathbb Q[[t]]$ take the projective hyperelliptic family
\[
 y^2=(x^2+1)^2(x^2-2)+t=x^6-3x^2-2+t.                       \tag{F.19}
\]
Its generic sextic has discriminant
$-46656\,t^2(t-2)(t-4)^2$, which is nonzero in
$\mathbb Q((t))$; the program checks this formula. Its special
fibre is integral with precisely the two ordinary nodes
$(x,y)=(i,0),(-i,0)$. The coefficient of $(x-i)^2$ in the
sextic at $i$ is $12\ne0$, and the same holds at $-i$.
At infinity the two points are smooth. Normalizing by
$w=y/(x^2+1)$ gives $w^2=x^2-2$, a smooth genus-zero curve;
the four inverse images of the nodes make this an irreducible
stable genus-two fibre. The total space is regular at each node
because the derivative of the defining equation with respect to
$t$ is $-1$. Thus this is a stable degeneration.
An element of $G_{\mathbb Q}$ taking $i$ to $-i$ interchanges
its two geometric nodes. By the graphic comparison (F.10), every
residue splitting gives this same permutation of their group
classes. Proposition 1.2 of \cite{CF_MochizukiComb2007} makes
the two classes distinct. Hence neither class is invariant under
the whole $G_{\mathbb Q}$. A finite extension fixes both nodes
and their branches and recovers exactly the conclusion of F17.6.
\end{proof}

Combining F17.1--F17.8, in residue characteristic zero the justified
arrows are a)$\Leftrightarrow$b)$\Rightarrow$c)$\Rightarrow$d) and
a),b)$\Rightarrow$e). The converses c),d)$\Rightarrow$a), the
arrows c),d)$\Rightarrow$e), and e)$\Rightarrow$a) all have the
explicit counterexamples just given. Condition e), with the finite
extension permitted in the manuscript, is equivalent to potential
good reduction under the finitely generated residue-field hypotheses
of F17.7. The predicate in f), ``pas
trop loufoque (\`a d\'efinir convenablement)'', has no definition
in the manuscript, so it supplies no additional mathematical
condition to put in an equivalence table.
